% id: 0272
% book: 라이트쎈 공통수학2 · 03 원의 방정식
% section: 유형 03 지름의 양 끝 점이 주어진 원의 방정식
% figure: none
% answer: ④
% answer_source: 해설
% variant_level: 0 (원본 전사)
% 이 파일은 build.mjs 가 items.json 에서 만든다. 고칠 때는 items.json 을 고친다.
\smash{\raisebox{1.5mm}{\dmkichul{라이트쎈 공통수학2 0272번 · 43쪽}}}
직선~$x+y-12=0$이 $x$축, $y$축과 만나는 점을 각각 $\pt{P}$, $\pt{Q}$라 할 때, 두 점~$\pt{P}$, $\pt{Q}$를 지름의 양 끝 점으로 하는 원의 방정식은?
\begin{choicesv}{$(x-3)^2+(y-3)^2=18$}{$(x-3)^2+(y-3)^2=72$}{$(x-6)^2+(y-6)^2=18$}{$(x-6)^2+(y-6)^2=72$}{$(x-6)^2+(y-6)^2=81$}\end{choicesv}
